%% kapa: exam=1 round=33 year=2022 period=1 form=A subject=civil_law subject_ko=민법 seq=14 number=14
%% source: Q-Net 감정평가사 기출문제 (cst00309) artlSeq=5209997 file=raw/1008/5209997/2022년도 제33회 감정평가사 1차시험 1교시(민법,경제학,부동산학).pdf
%% conversion: kordoc 4.24.51 -> tools/gen1.py (auto); review=proofread-sonnet; flags=none
\begin{question}{14}
복대리에 관한 설명으로 옳지 않은 것은? (다툼이 있으면 판례에 따름)
\begin{choices}
  \item 대리권이 소멸하면 특별한 사정이 없는 한 복대리권도 소멸한다.
  \item 복대리인의 대리권은 대리인의 대리권의 범위보다 넓을 수 없다.
  \item 복대리인의 대리행위에 대해서는 표현대리가 성립할 수 없다.
  \item 법정대리인은 그 책임으로 복대리인을 선임할 수 있다.
  \item 임의대리인은 본인의 승낙이 있거나 부득이한 사유있는 때가 아니면 복대리인을 선임하지 못한다.
\end{choices}
\end{question}
\answer{3}
